\chapter{Où le trouver, et comment l'installer}

Jenga s'obtient de trois façons, et elles ne servent pas le même usage. Ce
chapitre les donne toutes les trois, avec la façon de vérifier laquelle vous
avez réellement --- ce qui n'est pas la même question.

\section{Ce qu'il faut avoir avant}

\begin{center}
\begin{tabular}{@{}lll@{}}
\toprule
\textbf{Il faut} & \textbf{Version} & \textbf{Pourquoi} \\
\midrule
Python & $\geq$ 3.8 & Jenga est écrit en Python \\
un compilateur & clang, gcc ou MSVC & Jenga l'appelle, il ne le remplace pas \\
\texttt{watchdog} & $\geq$ 2.1 & seulement pour \texttt{jenga watch} \\
\texttt{colorama} & $\geq$ 0.4.4 & couleurs sous Windows \\
\bottomrule
\end{tabular}
\end{center}

\begin{nkretenir}
\textbf{Les deux dernières lignes sont facultatives}, et le fichier de
dépendances le dit lui-même : « optional but recommended for 'watch' command ».
Jenga construit sans elles.
\end{nkretenir}

\section{Voie 1 --- le paquet installé}

C'est la voie normale. Jenga se distribue comme une roue Python
(\emph{wheel}) ; l'installer place une commande \texttt{jenga} dans votre
\texttt{PATH} :

\begin{nkcode}{pyproject.toml:37-38}
[project.scripts]
jenga = "Jenga.Jenga:main"
\end{nkcode}

\begin{nkcode}{Installer}
pip install jenga-2.4.0-py3-none-any.whl --force-reinstall
\end{nkcode}

\begin{nkretenir}
\textbf{\texttt{-{}-force-reinstall} n'est pas une précaution superstitieuse.}
Sans lui, \texttt{pip} peut considérer qu'une version portant le même numéro est
déjà là et ne rien faire --- en affichant un message de succès. Vous croiriez
avoir installé la neuve.
\end{nkretenir}

\section{Voie 2 --- depuis les sources, en mode éditable}

Pour travailler \emph{sur} Jenga, on l'installe en mode éditable : la commande
pointe alors vers l'arbre de sources, et une modification prend effet sans
réinstaller.

\begin{nkcode}{Dans le depot de Jenga}
python -m pip install -e . --no-build-isolation
\end{nkcode}

Le dépôt fournit un script qui fait tout : nettoyage, réinstallation éditable,
construction de la roue et des archives d'exemples.

\begin{nkcode}{Dans le depot de Jenga}
./cri.sh              # dev + distribution
./cri.sh --tests      # ... et la suite pytest
\end{nkcode}

\begin{nkpiege}
\textbf{Ce script est destructeur, et dans cet ordre :} il fait
\texttt{rm -rf build dist} \emph{puis} \texttt{pip uninstall Jenga}, et
\emph{seulement ensuite} réinstalle. Un échec en cours de route vous laisse sans
commande \texttt{jenga}.

\textbf{La parade tient en une ligne, et elle se pose AVANT} : copiez
\texttt{dist/*.whl} ailleurs. La panne devient réparable par
\texttt{pip install <roue sauvegardée> --force-reinstall}. Un filet posé après
l'accident n'est pas un filet.
\end{nkpiege}

\section{Voie 3 --- embarqué dans NKCode}

\textbf{NKCode contient sa propre copie de Jenga.} C'est ce qui rend les boutons
\emph{Construire} et \emph{Exécuter} de l'éditeur autonomes : ils ne dépendent
pas de ce que vous avez installé.

\begin{nkretenir}
Le mécanisme est explicite dans le descripteur de NKCode : il \textbf{extrait la
roue officielle} --- environ 1 Mo --- dans \texttt{Build/\_jenga-embed/}, et
embarque ce contenu.

\begin{nkcode}{Applications/NKCode/NKCode.jenga}
stage = _norm(os.path.join(_REPO_ROOT, "Build", "_jenga-embed"))
# Re-extrait a chaque evaluation : un wheel plus recent doit gagner,
# sinon on embarquerait silencieusement une version perimee.
\end{nkcode}

Il cherche la roue dans \texttt{dist/} du dépôt Jenga, en remontant plusieurs
niveaux, et la variable d'environnement \texttt{NKCODE\_JENGA\_SRC} permet de
lui indiquer un autre emplacement.
\end{nkretenir}

\begin{nkpiege}[Deux versions dans la même session]
NKCode peut construire avec une version et votre terminal avec une autre. Vous
auriez alors \textbf{deux systèmes de construction silencieusement différents},
et un défaut d'outil se chercherait dans le code.

Le descripteur compare donc les deux et le \emph{dit} :

\begin{nkcode}{Applications/NKCode/NKCode.jenga}
[NKCode] Jenga embarque depuis le wheel : jenga-2.4.0-py3-none-any.whl
[NKCode] ATTENTION : le wheel embarque est en 2.2.0 alors que le jenga du
[NKCode]   systeme est en 2.4.0. NKCode construira donc AUTREMENT que votre
[NKCode]   terminal. Remede : ./cri.sh dans le depot Jenga.
\end{nkcode}
\end{nkpiege}

\begin{nkretenir}[Le défaut qui rendait cette comparaison muette]
Ce tri de roues se faisait autrefois sur le \emph{nom}, donc en texte. En ordre
lexicographique, \textbf{\texttt{2.9.0} est supérieur à \texttt{2.10.0}} : le
jour du passage à deux chiffres, NKCode aurait choisi la plus ancienne, en
silence.

Il compare désormais des \emph{tuples d'entiers}. \emph{Comparer des versions en
texte marche jusqu'à 9 et casse à 10} --- et la panne se présente comme « l'outil
annonce la mauvaise version », à des lieues de sa cause.
\end{nkretenir}

\section{Vérifier ce qu'on a vraiment}

\begin{nkcode}{Capture reelle}
$ jenga --version

             Multi-platform C/C++ Build System v2.4.0

Jenga version 2.4.0
\end{nkcode}

\begin{nkpiege}[Le sujet mesuré n'est pas toujours celui qu'on croit]
\textbf{\texttt{pip show jenga} et \texttt{jenga -{}-version} peuvent ne pas
parler du même programme.} Un poste peut porter plusieurs Python : \texttt{pip}
appartient à l'un, et le lanceur \texttt{jenga} de votre \texttt{PATH} peut venir
des \texttt{Scripts} d'un autre.

Mesuré sur le poste de développement : \texttt{python} résout vers un Python
3.14, tandis que \texttt{jenga} vient des \texttt{Scripts} d'un Python 3.13.

\textbf{Le sujet à mesurer est \texttt{jenga -{}-version}} --- la commande que vous
tapez --- jamais \texttt{pip show} seul.
\end{nkpiege}

\section{Le module d'aide à l'édition}

Vous verrez apparaître un dossier \texttt{.jenga-typings/} à côté de vos
descripteurs. Il contient deux fichiers, et le premier avoue sa nature :

\begin{nkcode}{.jenga-typings/jengaconfig.py}
# Module no-op genere par Jenga. `from jengaconfig import *` est resolu
# par jengaconfig.pyi cote editeur ; au runtime les symboles viennent de
# la propagation useconfig(). Ne rien mettre ici.
\end{nkcode}

\begin{nkretenir}
\textbf{\texttt{from jengaconfig import *} n'importe rien.} C'est une ligne
\emph{pour l'éditeur} : elle donne l'auto-complétion via le fichier
\texttt{.pyi}. À l'exécution, les symboles arrivent par la propagation de
\texttt{useconfig()}.

C'est un cas rare et honnête d'une ligne qui ne sert qu'à l'outillage --- et
qu'il faut savoir reconnaître, sans quoi on la cherche dans le mécanisme.
\end{nkretenir}

\begin{nkbilan}
Trois voies : la roue installée pour s'en servir, le mode éditable pour
travailler dessus, et la copie embarquée dans NKCode pour que l'éditeur soit
autonome. La troisième peut diverger des deux autres, et NKCode le signale ---
parce que deux systèmes de construction dans une même session font chercher un
défaut d'outil dans le code. Vérifiez toujours par \texttt{jenga -{}-version}, qui
interroge la commande réelle, et non par \texttt{pip show}, qui peut parler d'un
autre Python. Enfin, \texttt{jengaconfig} ne sert qu'à votre éditeur.
\end{nkbilan}

\begin{nkquestions}
\begin{enumerate}
\item Quelles sont les trois façons d'avoir Jenga, et à quoi sert chacune ?
\item Pourquoi \texttt{-{}-force-reinstall} lors de l'installation d'une roue ?
\item Dans quel ordre \texttt{cri.sh} désinstalle-t-il et réinstalle-t-il, et
      qu'est-ce que cela implique ?
\item Pourquoi NKCode compare-t-il sa version embarquée à celle du système ?
\item Que fait \texttt{from jengaconfig import *} à l'exécution ?
\end{enumerate}
\end{nkquestions}

\begin{nkpratique}
Installez Jenga, puis établissez --- par des commandes, pas par déduction --- les
quatre faits suivants :

\begin{enumerate}
\item la version que \texttt{jenga} annonce ;
\item la version que \texttt{pip} croit avoir installée ;
\item le chemin exact du lanceur \texttt{jenga} (\texttt{where} ou
      \texttt{which}) ;
\item l'interpréteur Python auquel ce lanceur est rattaché.
\end{enumerate}

Si les points 1 et 2 diffèrent, vous avez deux Python. Notez lequel gagne.
\end{nkpratique}

\begin{nkdemo}
Prouvez que le mode éditable est bien \emph{éditable}.

Installez Jenga par \texttt{pip install -e .}, puis ajoutez une ligne
\texttt{print("MARQUE")} au tout début de \texttt{Jenga/Jenga.py}. Relancez
\texttt{jenga -{}-version} : la marque doit apparaître, \textbf{sans réinstaller}.

Retirez la ligne, relancez : elle disparaît.

\textbf{Ce que la démonstration établit} : la commande lit l'arbre de sources à
chaque exécution. Refaites ensuite le même essai après une installation par
\emph{roue} --- la marque n'apparaîtra pas, parce que le code exécuté est une
copie. C'est la paire qui prouve, pas le premier essai seul.
\end{nkdemo}

\begin{nklogique}
Vous corrigez un défaut dans Jenga, vous relancez votre construction, et rien ne
change.

\begin{enumerate}
\item Énumérez au moins quatre causes possibles, de la plus banale à la plus
      retorse.
\item Pour chacune, donnez une commande dont le résultat \emph{diffère} selon
      que la cause est celle-là ou une autre.
\item L'une de ces causes ne se voit que si l'on regarde \emph{quel} programme a
      construit --- laquelle, et pourquoi les autres vérifications passent-elles
      à côté ?
\end{enumerate}
\end{nklogique}
